\gddchapter{Interview sequence}


Below is the guideline for the semi-open guided interview. 
The interview starts with easy warm-up questions and progresses towards harder questions as it goes along. 

The answers were translated from polish back to english.

\section{Introduction}
Remind the participant that this is not a test, rather a joint exploration.

\section{Icebreakers}

\begin{QandA}
   \item Now that you’ve tried both versions, how are you feeling? Was there anything that immediately surprised you?
         \begin{answered}
            Zuzia: No, I wouldn't say there was anything that surprised me.
            I would say I prefer the first version because there wasn't the delay between saying something and clicking it.
            But it was also a nice experience that it was working after the case.
            Yeah, because normally, for me, it doesn't work.
            So that was a nice experience.
            I would say, I don't really know how to say it in English.


            Gabriel: The way, okay, the ability, let's say.


            Zuzia: Yeah, the ability to make things do a work and they have more impact.
            Maybe translation the other way, yeah.
            \end{answered}

   \item If you had to describe the difference between these two games to a person who hasn't played either, what would you say?
         \begin{answered}
            That it would be more tiring the second version after a longer time because you can't just sit and click you can't sit and click on things you have to say them so I guess it's more tiring after some time but I would imagine for some people
            Problems with accessibility that would be a nice alternative I would say, you know to click on things and for example If there will be an audio description with it, so it would be you know more more accessible or even not more but They could be able to play if you're applying so that would be nice But I would say
            For me, as a person who likes to sit quiet in her own time, the first alternative would be nicer.
         \end{answered}
\end{QandA}




\section{Grand tour questions}
(remember to pause here 3-5 seconds)

\begin{QandA}


   \item From the moment you first used your voice to navigate, what were you thinking as you decided what to say to the system?

         \begin{answered}
            I have background with computers, so my first thing, and also with AI.
            and stuff like that, so my first thing was to say something that would be the most understandable for a computer.
            And then to see how much the range of things I can say that he will understand.
            You know, like, we can add guns and see if we can skip two places to go back to the first one.
         \end{answered}

    \item How was it like to use the voice input in the game?
         \begin{answered}
            I mean, it was nice.
            It understood me.
            And it was sometimes funny, like, you can't go back to lunch from here.
            It would be nice to go to lunch, I guess, in the situation that they have If they had to go for a food run, so it would be nice to go back to lunch, not lounge I would say I mean, I like trying new things, so it was nice for me to experiment with it, like, can I pick hands and everything, how much you can do by clicking and how much the computer, when you say things, processes
            So yeah, that was fun in this way.
            But still, I'm not the biggest fan of speaking when playing games.
            For example, there is a game that is called Yap Yap, and it's about mages right now.
            And you cast spells by saying them.
            This one was much better, the accessibility of voice than the one in the game, because I have a problem with casting spells.
            So I turned it off and turned
            option with clicking buttons and it's much more enjoyable for me so I guess it's me but this one was more pleasant than the the other one because you know sometimes polish words were taken for other things and I think if this one understood much better because there you have to say like aero for from you know aerodynamic and everything but it only works when you say euro
            So I guess lunch is closer than euro and euro.
         \end{answered}


\end{QandA}




\section{Core comparative questions}

\begin{QandA}
   \item In which version did you feel more 'inside' the story world?
         \begin{answered}
            Zuzia: The first one.
            can think in my mind I go there but when I'm saying go back there as a immersive play style you don't really speak you go there you know you don't say to your legs go there maybe it would be better if it was in the first person maybe if I would be speaking in the first person but my mind was like you know it was written like that so I'm not saying
            will go there now I want to pick up maybe this way would be more immersive but as I did it the first one is more immersive because it happens in my head and doesn't really go out in the realistic world that is my couch right now you know I understand I understand of course of


            Gabriel: have some experience with playing TRPG games like Dungeons \& Dragons and stuff like that and when you control your character do you control it with the first person or third person?


            Zuzia: Depends because sometimes when I explain things to people I say you see now that the character does this and this and the spell looks like that or like that yeah but I say I want to go there and do that
            It depends on the narrative, if I'm explaining things to someone, what they see, to make them more immersed or outplaying me.
            It's hard for me sometimes because it's another wall that you have to cross when you say I go there or I say that and you say it and don't explain what you are saying but mostly in the first person.


            Gabriel: What would you think would be the difference that makes you use second person when you played D \& D compared to the fact that you chose the third person.


            Zuzia: I didn't even chose the third person.
            I was speaking in the first person to the computer.
            That's the reason.
            Because I was saying go to launch.
            So I'm not saying I want to go to, character goes to, but I'm speaking to the computer that he has to get me there.
            He has to go there, not me.


            Gabriel: okay because when you control your character you are in character and myself and yes there is not not a layer of control that you have to go through there is not a dungeon master that says i will take you there right i'm not speaking to my friend go to


            Zuzia: yeah this place so i can do that and that i'm saying i want to go there maybe it's because there is more people and he's just a narrator but for me i guess computer is a tool and also when we speak about siri and we were talking about siri our siri is also a tool that you use i see i'm trying to understand if there's like a


            Gabriel: authority here like when you play D\&D you are the ultimate authority that decides what can happen outside of the dungeon master telling you that you cannot do something but you say that I'm going there I'm doing that and yeah and you're not you you're not piloting anything yeah okay okay and speaking of the effort also how did the effort of thinking what to say compare to the effort of knowing which key to press


           
         \end{answered}

   \item How did the 'effort' of thinking of what to say naturally compare to the 'effort' of knowing which keys to press?

         \begin{answered}
            the only thing that was different is that you didn't see the comments you have to think them in their head and I had to do the rerun from the start but because I have a good memory I remembered the exact places it could be more immersive because you have to think what can you do but again in this version you don't have a
            many opportunities to do things because you can't go somewhere else that is not prescribed in the exact room.
            You know, you can go back two times, you can only go one time back or forward.
            So I think you have to think like the comments you had before.
            So you automatically thinking not out of the box.
            So there isn't much effort, but
            It's also, I think, because I'm...
            Sorry, he has a strange [?].
            It's also because I know the comments and I'm thinking more computerically.
            Yeah, I'm thinking more like a computer than I think many other persons.
            I'm going to start thinking how it was presented to me at first.
            
            But I think it could be frustrating sometimes if I didn't see the numbers version before because I wouldn't know how the program is thinking and what should I say and how it would work.
         \end{answered}

    \item Did the freedom to say 'anything' make you feel more in control, or was it more overwhelming than having a fixed list of keys?

         \begin{answered}
            Zuzia: no difference for me okay but i have this problem when something is presented for me before it very gets in my mind the same i have i think because of the spectrum that i'm on you know when i have rules that are shown to me i don't really look for ways to get around them i really stick to them from the start to the end
            So I have some problems in life because of that.
            Because I'm thinking less out of the box.
            So for me it's natural and I don't really think in other ways.
            I don't have to feel the other aspects of it that you are speaking of.
            And it was, yeah, the control, the sprawczość that I was talking about.
            It was nice.
            And it was nice that I can say something and it's processing.
            And then it works.


            Gabriel: That felt good in a way.


            Zuzia: Yeah, because it's something new.
            I never had this experience before.
            Even with people, I guess.
         \end{answered}
    \item How would you compare the overall experience of using the two game versions? Which one would you choose and why?

         \begin{answered}
            For a moment I would say the second version, because it's something new, but long time you have to play.
            For some hours I would choose the first one, because it's just more immersive and easier in the long term to use.
            You know, you just click the buttons and it's not like that you're sitting alone in a computer and you're just saying, go to launch, pick up the crowbar.
            No, I'm maybe not the most expressive person when it comes to comments because I'm just built like that.
            But I would feel more comfortable with the first version.
         \end{answered}
\end{QandA}






\section{Probing}
In this part of the interview look into the sticky moments noted above and ask follow-up questions probing for more detail.
Include lexical reflection (eg. you used the word "clunky" when explaining navigation. Could you tell me more about what was happening
in that moment?)

Most of the probing in this interview happened as a follow up to the question "In which version did you feel more 'inside' the story world?"


\section{Final reflection}

\begin{QandA}
   \item After reflecting on everything today, is there anything else you’d like to add about using your voice to play this game that we haven't talked about?

         \begin{answered}
            Zuzia: You know, if I didn't have the time, I would probably read more.
            of the text more immersively, not just go through this because I wanted to see the most of the things that you came up with, yeah.
            Most rooms and everything like that.
            And I was interested if there will be a plot twist and more dense action than going through places.
            And also because you have things like, you know, you have the
            Keyhole that was scratched and someone tried to open it in a korean store You know Then you have metal key card And okay, it doesn't match there, but for example why I can't use a crowbar If I could use a crowbar before Because I could do that if I wanted So it's better not to include things like that
            in the same place if I can use the crowbar maybe if there is no use for a crowbar later you can just discard it after the opening the clothing store doors because then you have a feeling that you could do it in the real life but the game doesn't let you I guess maybe in the supermarket store or later you could find the key for the doors in the korean store but still you
            it was a little bit not satisfying that I couldn't use it.


            Gabriel: It takes away your agency in a way, right?


            Zuzia: Yeah.
            It's better not to include it there.
            Or maybe if you go back to the room, maybe it would be more hard because, you know, if you would want to go back after, then you can say, look for, you know, get a new option in game, look for the way out or something.
            But another thing, if you are
            using voice control and you see that you don't have that option there it wouldn't work, you know because with voice option I would like to try to use the crowbar because I would see I wouldn't see the comments so I would not know what are my limits and I would try to use things and try for example to say different words so it would work but it wouldn't work so I would get frustrated I guess
            You know, the voice comment is nice, but also if you don't know your limitation, it's easy to get frustrated because you want to say stuff and you don't know if there are keywords that will work.
          
         \end{answered}
\end{QandA}

\section{Closing}
Thank the participant for his time and tell them again how much value his input and time gives me.






